% !TeX program = pdflatex
\documentclass[10pt,letterpaper]{article}
% Compiler Engineering Suite - shared standalone pdfLaTeX preamble
\documentclass[10pt,letterpaper]{article}
\usepackage[T1]{fontenc}
\usepackage[utf8]{inputenc}
\usepackage{lmodern}
\usepackage[margin=.78in,headheight=18pt]{geometry}
\usepackage[table]{xcolor}
\usepackage{microtype}
\usepackage{amsmath,amssymb}
\usepackage{booktabs,longtable,array}
\usepackage{enumitem}
\usepackage{fancyhdr}
\usepackage{hyperref}
\usepackage{needspace}
\usepackage{fancyvrb}
\usepackage{tikz}
\usetikzlibrary{arrows.meta,positioning,shapes.geometric}
\definecolor{CompilerBlue}{HTML}{1F4E79}
\definecolor{CompilerInk}{HTML}{263746}
\definecolor{CompilerGray}{HTML}{52616B}
\definecolor{CompilerLight}{HTML}{EFF4F8}
\hypersetup{colorlinks=true,linkcolor=CompilerBlue,urlcolor=CompilerBlue,citecolor=CompilerBlue}
\setlength{\parindent}{0pt}
\setlength{\parskip}{6pt plus 2pt}
\setlist[itemize]{leftmargin=*,itemsep=3pt,topsep=2pt}
\setlist[enumerate]{leftmargin=*,itemsep=3pt,topsep=2pt}
\pagestyle{fancy}
\fancyhf{}
\fancyhead[L]{\small\textcolor{CompilerGray}{Compiler Engineering Documentation Suite}}
\fancyhead[R]{\small\thepage}
\renewcommand{\headrulewidth}{.4pt}
\newcommand{\evidence}[1]{\par\smallskip\noindent\textbf{Evidence required:} #1\par}
\newcommand{\sourceanchor}[1]{\textbf{Book anchor:} #1\par}
\newcommand{\modern}{\textbf{Modern engineering extension (not attributed to the source book):} }
\newcommand{\writinglines}[1]{\par\noindent\rule{\linewidth}{.3pt}\par\vspace{#1\baselineskip}\noindent\rule{\linewidth}{.3pt}\par}
\newcommand{\checkboxline}[1]{\noindent $\square$\quad #1\par}
\newcommand{\formfield}[2]{\noindent\textbf{#1}\par\vspace{#2\baselineskip}\hrule\medskip}

\hypersetup{pdftitle={Source Crosswalk and Provenance Register},pdfauthor={Compiler Engineering Documentation Suite}}
\begin{document}
\begin{titlepage}
\centering\vspace*{1.0in}
{\Huge\bfseries\textcolor{CompilerBlue}{Source Crosswalk and Provenance Register}\par}
\vspace{.22in}{\Large Fourteen chapters, appendices and suite artifact references\par}
\vspace{.4in}{\large Compiler Engineering Documentation Suite\par}
\vfill
\begin{minipage}{.86\linewidth}\small
Source basis: Keith D. Cooper and Linda Torczon, \emph{Engineering a Compiler}, Third Edition (Morgan Kaufmann/Elsevier, 2023). All explanatory text, examples, questions, and implementation designs in this document are independently authored. The book is not included.\par\medskip
Document class: Technical reference.
\end{minipage}
\vfill\textcolor{CompilerGray}{Version 1.0 --- 2026-10-30}\par
\end{titlepage}
\tableofcontents\clearpage
\section{Chapter and appendix source register}
\begin{longtable}{|p{0.11\linewidth}|p{0.37\linewidth}|p{0.45\linewidth}|}
\hline
Chapter & Canonical heading & Engineering deliverable\\ \hline\endfirsthead
\hline
Chapter & Canonical heading & Engineering deliverable\\ \hline\endhead
1 & Overview of Compilation & Pipeline diagram\\ \hline
2 & Scanners & State-transition drawings\\ \hline
3 & Parsers & Grammar proof\\ \hline
4 & Intermediate Representations & IR schema\\ \hline
5 & Syntax-Driven Translation & Typed AST\\ \hline
6 & Implementing Procedures & Stack map\\ \hline
7 & Code Shape & Code-shape inventory\\ \hline
8 & Introduction to Optimization & Before/after IR\\ \hline
9 & Data-Flow Analysis & Iteration log\\ \hline
10 & Scalar Optimization & Legality matrix\\ \hline
11 & Instruction Selection & Selection table\\ \hline
12 & Instruction Scheduling & DAG\\ \hline
13 & Register Allocation & Graph\\ \hline
14 & Runtime Optimization & State machine\\ \hline
\end{longtable}

\section{Cross-suite dependencies}
\begin{itemize}
\item Parsing Techniques: reuse formal grammar prerequisites but keep compiler lowering and optimization evidence here.
\item Algorithm Engineering: reuse graph/fixed-point analysis and complexity techniques; do not duplicate book-derived explanations.
\item Existing \texttt{src/programming/technical-documents/engineering-compiler-algorithms-pseudocode.tex}: related preserved pseudocode reference; this suite adds exercises and executable tests, not a replacement.
\item Language-specific secure coding: cross-reference C/C++ and Python guidance without relocating language-specific standards into this suite.
\end{itemize}

\section{Citable source metadata}
Keith D. Cooper and Linda Torczon. \emph{Engineering a Compiler}, Third Edition, Morgan Kaufmann, 2023, ISBN 978-0-12-815412-0. The chapter and appendix section ranges come from the uploaded edition\textquotesingle s table of contents. The suite\textquotesingle s tests, examples and security materials are independently created.

\end{document}
